\documentclass[11pt]{article}
\usepackage[margin=1.15in]{geometry}
\usepackage{amsmath,amssymb,amsthm}
\usepackage[hidelinks]{hyperref}
\newtheorem{theorem}{Theorem}
\newcommand{\E}{\mathbb E^t_{x,x}}
\title{Gauss--Bonnet via Brownian Bridges}
\author{}
\date{}
\begin{document}
\maketitle
\begin{abstract}
The McKean--Singer formula expresses the Euler characteristic as a heat
supertrace. On an oriented surface, its Brownian bridge representation
separates a scalar curvature weight from the trace of stochastic parallel
transport. The latter differs from the identity trace by order $t^2$ in
expectation, leaving the Gaussian curvature as the leading term.
\end{abstract}

\begin{theorem}
Let $(\Sigma,g)$ be a closed oriented Riemannian surface. If $K$ is its
Gaussian curvature and $dA$ its area measure, then
\[
  \frac{1}{2\pi}\int_\Sigma K\,dA=\chi(\Sigma).
\]
\end{theorem}

\section{The heat supertrace}
Set $D=d+d^*$ and $\Delta_k=(dd^*+d^*d)|_{\Omega^k}$. Hodge theory
identifies $\ker\Delta_k$ with $H^k_{\mathrm{dR}}(\Sigma)$. For every
positive eigenvalue $\lambda$ of $D^2$, the restriction of $D$ maps the
even eigenspace isomorphically to the odd eigenspace, with inverse given
by the opposite restriction of $\lambda^{-1}D$. The positive spectrum
therefore cancels in the supertrace. Writing $H_k$ for the heat kernel
of $\Delta_k$, this gives the McKean--Singer identity
\[
  \chi(\Sigma)=\int_\Sigma
  \bigl(2p_t(x,x)-\operatorname{tr}H_1(t,x,x)\bigr)\,dA(x),
  \qquad p_t=H_0 .
\]
Here the Hodge star identifies the heat kernels in degrees zero and
two. This identity holds for every $t>0$.

\section{The pinned path representation}
Choose Brownian motion with generator $-\Delta_0$, so $p_t$ is its
transition density and $p_t(x,x)\sim (4\pi t)^{-1}$ uniformly. With
generator $-\Delta_0/2$ instead, the same heat operator uses a bridge of
duration $2t$, and the potential in its exponential is $K/2$.

Under the bridge from $x$ to $x$ of duration $t$, let
$A_t=\int_0^t K(B_s)\,ds$ and let $U_t$ be stochastic parallel transport
on $T_x^*\Sigma$ along the bridge. The Weitzenb\"ock identity on one-forms
is $\Delta_1=\nabla^*\nabla+K\,\mathrm{Id}$. The covariant Feynman--Kac
formula in pinned form \cite[Sec.~4]{hsu} consequently gives
\[
   H_1(t,x,x)=p_t(x,x)\,\E\bigl[e^{-A_t}U_t^{-1}\bigr].
\]
The factor $p_t(x,x)$ is the return density. Since the connection is
metric and orientation preserving, $U_t$ is a planar rotation with
angle $\theta_t$, and its inverse has the same trace $2\cos\theta_t$.
The pointwise integrand in the heat supertrace becomes
\begin{equation}\label{eq:integrand}
  2p_t(x,x)\left(1-\E[e^{-A_t}\cos\theta_t]\right).
\end{equation}

\section{The small-time limit}
The factor $p_t(x,x)$ grows like $t^{-1}$, so convergence of
$\cos\theta_t$ to $1$ alone is insufficient. Hsu's short-bridge estimates
\cite[Prop.~2 and Sec.~6]{hsu}, with time rescaled for our generator,
give uniformly in $x$
\[
   \E\|U_t-I\|_{\mathrm{HS}}^2=O(t^2),\qquad
   \E\sup_{0\le s\le t}d(B_s,x)=O(\sqrt t).
\]
The holonomy scale can be understood in a frame obtained by radial
parallel transport from $x$. The connection form vanishes at the
center, and the leading stochastic line integral has the scale of
stochastic area, namely $t$. For a planar rotation the trace depends
quadratically on this displacement, since
\[
    1-\cos\theta_t=\tfrac14\|U_t-I\|_{\mathrm{HS}}^2.
\]
As $|A_t|\le t\|K\|_\infty$, it follows that
$\E[e^{-A_t}(1-\cos\theta_t)]=O(t^2)$.

Smoothness and compactness make $K$ Lipschitz. The bridge displacement
bound thus gives $\E A_t=tK(x)+O(t^{3/2})$. Expanding the exponential
uniformly as $e^{-A_t}=1-A_t+O(t^2)$ yields
\[
  1-\E[e^{-A_t}\cos\theta_t]=tK(x)+O(t^{3/2}).
\]
The sign of the linear term in the exponential is negative. The
subtraction in the heat supertrace produces the positive curvature
term. Using the uniform scalar heat-kernel asymptotic in
\eqref{eq:integrand} gives the local limit $K(x)/(2\pi)$. Integration
over the compact surface proves the theorem. \hfill$\square$

\begin{thebibliography}{9}
\bibitem{hsu} E.~P.~Hsu, \emph{Stochastic local Gauss--Bonnet--Chern theorem},
Journal of Theoretical Probability \textbf{10} (1997), 819--834.
\href{https://sites.math.northwestern.edu/~ehsu/Stochastic\%20Local\%20Gauss-Bonnet-Chern\%20Theorem.pdf}{Author's copy}.
\end{thebibliography}
\end{document}

